\section{Related Work}
\label{app:related}

\paragraph{Automated research and idea generation.} Language models are
increasingly used to carry out parts of the research process itself.
\citet{wang2024scimon} retrieve inspirations from past literature and optimise
explicitly for novelty by comparing a candidate against prior papers;
\citet{baek2025researchagent} refine proposals against a set of LLM reviewing
agents that give structured feedback along several dimensions; and
\citet{lu2024aiscientist} and \citet{schmidgall2025agentlab} close the loop
from ideation through experiments to a
written paper and a simulated review. \citet{si2025llmideas} evaluate what such
pipelines produce, finding the generated ideas rated more novel than
human-written ones but less feasible; \citet{zheng2025automation} survey the
progression of these systems toward greater autonomy.
Every one of these systems contains an internal novelty check, and it is the
check rather than the generator that bounds them: once proposals are cheap, the
output of the pipeline is whatever its filter admits. The filter is nonetheless
the component these systems evaluate least directly, judged only through the
ideas that eventually emerge downstream of it.

\paragraph{Judging the novelty of scientific work.} Studied on its own, the
check has passed through several generations of method. The earliest operate on
citation statistics, measuring how unusual the combinations of references in a
paper are against historical co-occurrence \citep{uzzi2013atypical}; they
require the work to already exist and be cited, and reduce novelty to patterns
of co-occurrence rather than to what a paper says. Document-level semantic
representations replace surface statistics with learned relatedness, scoring how
close a submission is to prior papers through citation-informed embeddings
\citep{cohan2020specter,ostendorff2022scincl}, but they compare documents as
points and recover nothing about which content is shared. Most recent systems
retrieve related literature and hand the comparison to a language model.
\citet{liang2024llmfeedback} obtain structured feedback from a single prompt, of
which significance and novelty is one section, and report partial overlap with
human reviews; \citet{zhu2025deepreview} train a
reviewer model on a multi-stage process of structured analysis, retrieval, and
evidence-based argumentation; \citet{chang2025treereview} decompose the review
into a tree of questions resolved from leaf to root, and
\citet{jin2024agentreview} simulate the review process itself to expose how
much of the outcome is attributable to reviewer bias rather than to the paper;
\citet{afzal2026beyond} split the judgment into
claim extraction, retrieval and synthesis of related work, and structured
comparison, motivating the design with an analysis of human-written novelty
reviews; and \citet{dasilva2025graphmind} build an interactive system on the
same retrieve-then-compare pattern. \citet{schopf2026rinobench} supply the
evaluation counterpart and report that model ratings still diverge substantially
from expert ones even when the accompanying reasoning appears sound. Two
properties persist across these generations, and both are structural rather
than incidental. The evidence is retrieved by the system under evaluation, so
its verdicts are unattributable: an error cannot be assigned to the retrieval or
to the judgment, and the two therefore cannot be improved, or even measured,
separately. We remove this by conditioning on a given, temporally admissible
evidence set, which additionally turns the evidence into a variable that can be
intervened on. The second property is that the target is scored as an undivided
whole, against the closest prior work or against a retrieved set read as a
block. Such a comparison returns one number per target, and no post-processing
of that number can recover which part of the idea it is about; the information
is discarded before the score is formed, not lost afterwards.

\paragraph{Decomposed verification against evidence.} Saying which part remains
requires decomposing the target, and outside novelty assessment this has become
the standard recipe for factuality evaluation: break a text into smaller
propositions and check each against retrieved evidence.
\citet{min2023factscore} decompose a passage into atomic facts, verify each
against a knowledge source, and report the supported fraction;
\citet{wei2024safe} scale the same procedure with search as the evidence source;
and in the scientific domain the corresponding task checks a stated claim
against candidate abstracts while identifying the supporting rationales
\citep{wadden2020scifact,wadden2022multivers}, where the choice of evidence
source itself moves the verdict \citep{vladika2024knowledge}. Three properties of the recipe do
not survive the move to research ideas. The propositions are assumed to be
given, whereas an idea carries no stated claims, and the decomposition step is
not a neutral preprocessor: \citet{wanner2024decomposition} show that scores of
this family shift materially with the decomposition used. Support is treated as
a label, whereas a prior work can cover part of a unit and not the rest. And
each unit is verified independently before the results are averaged, so the
score is a function of the per-unit coverages alone and cannot express whether
a combination of units has precedent. We accordingly induce the units rather
than reading them off the text, measure support as a graded quantity, and add a
residual defined over pairs of units.
